\chapter{Reversing a finite surgery history}
\label{ch:finite-history-v4}

This chapter reverses one specified surgery flow after its extinction time has
been found.  Every factor below is an actual component supplied by the local
surgery theorem for this flow.  In particular, we never replace the discarded
components by a separately classified list.  The result is a finite smooth
connected-sum assembly of the actual time-zero slice.  The finite reversal is
the argument of \cite[Corollary 15.4, pp.~358--359]{MT2007}; its local input is
\cite[Proposition 15.3, pp.~357--358]{MT2007}.

Fix a normalized initial metric $g_0$ on a smooth three-manifold $M$.  The
dimension is three throughout. Fix the flow $F$ of
Section~\ref{sec:global-assembly}, and write $M_t$
for its slice at time $t$,
let $H$ be the extinction time supplied
by Theorem~\ref{prop:v4-global-extinction} for this same $F$, and assume
that $M$ is connected. The local surgery topology theorem,
Theorem~\ref{thm:surgery-local-topology}, supplies, for every surgery time
$T\leq H$, an event record whose incoming reference slice is an actual slice
$M_{q_T}$ with $0\leq q_T<T$.  We assume $M$ is nonempty, compact,
Hausdorff and second countable, with the smooth and measurable structures of
the global-flow input.  The chosen time satisfies
$H\in\operatorname{Surg}(F)$ and $M_H=\varnothing$.
The global flow is defined for all nonnegative times, and zero is not a
surgery time.

The chronology used by the reversal has the following form.  There is an
anchor $p_T<q_T$, with either $p_T=0$ or
$p_T\in\operatorname{Surg}(F)$, such that $[p_T,q_T]$ lies in the time
domain and there is no surgery in $(p_T,q_T]$ or $(q_T,T)$.
There is a specified diffeomorphism $M_{p_T}\to M_{q_T}$ equal to the
regular-slab transport.  All event conclusions and cap correspondences come
from the local topology theorem for this same $F$.  These hypotheses explain
which identifications, in addition to extinction, the reversal uses.
More precisely, $q_T$ lies after the lower endpoint of the actual event's
preterminal slab, both in the nonempty and vanishing branches.  The local
conclusion on $M_{q_T}$ is the conclusion for that actual event transported
along its preterminal identification.
We construct these reference times below from local finiteness and the
event's own regular interval.

\section{The event ledger and its factors}
\label{sec:finite-event-ledger}

The flow's local finiteness says that
\[
  \mathcal S_K=\{T\in\operatorname{Surg}(F):T\in K\}
\]
is finite for every compact time set $K$.  Applying it to the compact interval
$[0,H]$ gives the event ledger
\begin{equation}
  \mathcal E=\operatorname{Surg}(F)\cap[0,H].
  \label{eq:finite-ledger}
\end{equation}
An event index is a time $T\in\mathcal E$, ordered by the real
times.  Since the chosen extinction event belongs to $\mathcal E$, the ledger
is nonempty and has a first element.

For $T\in\mathcal E$, write
\[
  E_T=(q_T,S_T),
\]
where $q_T<T$ is the selected pre-event reference time and
\[
  S_T:\;M_{q_T}\longleftarrow M_T
\]
denotes the local topology conclusion, rather than a map between the two
slices.  Its finite index set is $J_T$; the carrier
of the $i$th local piece is $P_{T,i}$ and its kind is one of
\[
  \mathrm{survivor},\qquad \mathrm{sphere\ bundle},\qquad
  \mathrm{positive\ spaceform}.
\]
A sphere bundle here is a smooth, surjective map onto $S^1$ with smooth local
trivializations having fiber $S^2$; both orientable and nonorientable bundles
are allowed.  A positive spaceform carries a Riemannian metric and its
Levi-Civita connection with constant strictly positive sectional curvature.
For a piece of either kind the local conclusion supplies the corresponding
bundle or metric data, not just a name for its topological type.
For a survivor, $S_T$ gives a region $U_{T,i}\subset M_T$ and a smooth region
equivalence
\[
  P_{T,i}\simeq M_T|_{U_{T,i}}.
\]
The survivor regions are pairwise disjoint connected components and cover the
post-event slice.  Every local piece is nonempty, compact and connected.  If $M_T$ is
empty, the event record asserts that no $i\in J_T$ is a survivor.  This last
clause is what makes the terminal ledger event a genuine base case.
If $M_T$ is nonempty, every inserted cap region $C_{T,a}\subset M_T$ is
contained in some survivor region: for every cap index $a$ there is an
$i\in J_T$ of survivor kind with $C_{T,a}\subset U_{T,i}$.
The reverse induction uses the whole-slice assembly; this cap condition
retains the event provenance at its input.

Define the exact discarded index set and its carrier by
\[
  D_T=\{i\in J_T:\operatorname{kind}(T,i)\neq\mathrm{survivor}\},
  \qquad P(T,i)=P_{T,i}\quad(i\in D_T),
\]
and define the global summand type
\begin{equation}
  \mathcal D=\coprod_{T\in\mathcal E}D_T.
  \label{eq:global-discarded-index}
\end{equation}
Thus a summand remembers both its event and its local index.  The associated
family $P:\mathcal D\to\{\text{slice carriers}\}$ is the literal family
$P(T,i)=P_{T,i}$; no quotient identifies factors from different events.
For $e\in\mathcal E$, define the tail index set
\[
  \operatorname{Tail}(e)=\{(T,i)\in\mathcal D:e\leq T\}.
\]

\begin{definition}[Smooth finite connected-sum assembly]
\label{def:finite-smooth-assembly}
\lean{PoincareMT.SmoothFiniteConnectedSumAssembly}
For a finite family $P_a$ and a carrier $X$, an assembly consists of an initial
carrier $C$, a partition $C=\coprod_a V_a$ into open-and-closed regions,
and diffeomorphisms $P_a\to V_a$, together with a reflexive-transitive chain
of smooth connected-sum steps from $C$ to $X$.  A step first decomposes its
source into two open-and-closed smooth carriers $B_1\sqcup B_2$.  On each side
it chooses an embedded ball parametrized on the radius-two ball in
$\mathbb R^3$, removes the closed radius-one ball, and identifies the two
punctured sides with disjoint open regions of its target.  A diffeomorphism
$\phi:S^2\to S^2$ and a smooth collar $S^2\times(-1,1)$ specify the joining
sphere.  In the negative collar half the map agrees with the first ball
coordinate at $(1-s)z$; in the positive half it agrees with the second ball
coordinate at $(1+s)\phi(z)$.  The punctured regions and central sphere cover
the target, with the central sphere disjoint from both regions.  All region
and collar maps have smooth inverses on their images.  The reflexive case is allowed, so an empty chain
of operations is an assembly as well.
\end{definition}

Thus an assembly retains the geometric operations as well as the list of
factors.  This is the object consumed by the later sphere-reduction chapter.
The two sides of an operation, and its target, need not be connected: the
operation can join two components while retaining all other components.
This flexibility is necessary because the intermediate surgery slices can
be disconnected.

\begin{lemma}[Transporting an assembly]
\label{lem:finite-assembly-transport}
\uses{def:finite-smooth-assembly}
\lean{PoincareMT.SmoothFiniteConnectedSumAssembly.nonempty_transportTarget}
If a family has an assembly with target $X$ and $d:X\to Y$ is a
diffeomorphism, the same family has an assembly with target $Y$.
\end{lemma}

\begin{proof}
If the assembly has no operation, its target is its initial disjoint union.
Use $Y$ as the new initial manifold, replace each region $V_a$ by $d(V_a)$,
and postcompose its identifying map with $d$.  The regions still form an
open-and-closed partition, and the new assembly again has no operation.
Otherwise retain every operation before the last one.  In the last
operation postcompose both punctured-region maps and its collar with $d$;
precompose their inverses with $d^{-1}$.  The two balls and the sphere
gluing map are unchanged.  Disjointness and coverage pass through $d$, and
the two collar identities pass through postcomposition.  Thus the last
operation now ends at $Y$.  In neither case is a diffeomorphism counted as
an additional connected-sum operation.
\end{proof}

\begin{lemma}[Choosing compatible event references]
\label{lem:finite-event-references}
\uses{lem:finite-assembly-transport,thm:surgery-local-topology}
\lean{PoincareMT.m72LocalTopologyFromRaw}
The local topology conclusions of the fixed flow can be placed at reference
times $q_T$ with all the predecessor and regular-interval properties stated
above.  This changes none of their pieces, survivor identifications, or
cap correspondences.
\end{lemma}

\begin{proof}
Fix an event $T$ and let $r_T<T$ be the lower endpoint of its stored
preterminal interval.  Local finiteness makes
\[
 B_T=\{0\}\cup\bigl(\operatorname{Surg}(F)\cap[0,T)\bigr)
\]
a nonempty finite set.  Put $p_T=\max B_T$.  Since $T>0$, every element
of $B_T$ is strictly less than $T$, so $0\leq p_T<T$.  Any surgery time
strictly between $p_T$ and $T$ would belong to $B_T$ and exceed its maximum.
There are consequently no such surgery times.

Choose
\[
 q_T=\frac{\max\{p_T,r_T\}+T}{2}.
\]
Both entries in the maximum are less than $T$.  Therefore
$p_T<q_T<T$ and $r_T<q_T<T$.  The interval $[p_T,q_T]$ is in the
nonnegative time domain, and the absence of surgery in $(p_T,T)$ gives
the two required surgery-free intervals.  The regular flow identification
on $[p_T,q_T]$ supplies the specified map $M_{p_T}\to M_{q_T}$.

The event's preterminal identification also gives a diffeomorphism
$M_{r_T}\to M_{q_T}$.  Apply
Lemma~\ref{lem:finite-assembly-transport} to the local reconstruction
whose target is $M_{r_T}$.  Only that target is changed: every local
piece, its kind, and its map into the outgoing slice remain the same.
For a nonempty outgoing slice this retains the actual cap-to-survivor
correspondence.  For an empty outgoing slice the vanishing-event
conclusion supplies the no-survivor rule, which is likewise unchanged.
Thus the two event branches supply the claimed data for the same $F$.
\end{proof}

\section{Successors and unchanged intervals}
\label{sec:finite-successors}

Let $T\in\mathcal E$ with $T<H$.  In particular this holds whenever an
$i\in J_T$ is a survivor: otherwise $T=H$ and the empty-event rule excludes
that kind.  The finite set
\[
  \{T'\in\mathcal E:T<T'\}
\]
is nonempty since it contains $H$.  Let $T'$ be its least element.  It is
minimal among all surgery times greater than $T$: times at most $H$ lie in
the ledger, and times greater than $H$ are also greater than $T'$.  The predecessor anchor in the
event record of $T'$ cannot be a time strictly between $T$ and $T'$: such a
time would itself be a surgery event and contradict minimality.  Hence
$p_{T'}\leq T$.  If $p_{T'}<T$, then either $T\leq q_{T'}$, contradicting
the absence of surgery in $(p_{T'},q_{T'}]$, or $q_{T'}<T<T'$, contradicting
the absence of surgery in $(q_{T'},T')$.  Thus $p_{T'}=T$.  In particular,
\begin{equation}
  T<q_{T'},\qquad [T,q_{T'}]\subset\operatorname{Dom}(F),
  \qquad \operatorname{Surg}(F)\cap(T,q_{T'}]=\varnothing.
  \label{eq:successor-slab}
\end{equation}
The regular-slab theorem supplies a diffeomorphism
\begin{equation}
  d_{T,T'}:M_T\longrightarrow M_{q_{T'}}.
  \label{eq:successor-diffeo}
\end{equation}
For a survivor $i$, write $U_{T,i}=\operatorname{Comp}_{M_T}(y)$ using the
local conclusion and put $x_{T,i}=d_{T,T'}(y)$.  A homeomorphism and its inverse
carry connected sets to connected sets; maximality of connected components
therefore gives
$x_{T,i}\in M_{q_{T'}}$ and
\begin{equation}
  d_{T,T'}(U_{T,i})=\operatorname{Comp}_{M_{q_{T'}}}(x_{T,i}).
  \label{eq:successor-component}
\end{equation}
Postcomposing the local region equivalence with $d_{T,T'}$ gives the required
componentwise equivalence from $P_{T,i}$ to this connected component.  This is
the successor record attached to $(T,i)$; its target event is again an element
of the same finite ledger because $T'\in[0,H]$.

The order and minimality assertions are used later, not merely recorded as
metadata.  They ensure that a survivor is transported across one surgery-free
interval, rather than across an interval containing an unaccounted event.

\section{The local substitution construction}
\label{sec:finite-substitution}

Suppose an assembly $R$ has already been constructed for the tail beginning
at a later event $T'$, with target $M_{q_{T'}}$.  Transporting the target of
$R$ backwards by $d_{T,T'}^{-1}$ using
Lemma~\ref{lem:finite-assembly-transport} gives an assembly of the whole slice $M_T$.
Let $C_T$ be the initial carrier of the local assembly, with its partition
$C_T=\coprod_{i\in J_T}V_{T,i}$ and identifications
$a_i:P_{T,i}\to V_{T,i}$.  Define the discarded carrier
$Z_T=\bigcup_{i\in D_T}V_{T,i}$ with its induced smooth structure.
The restricted partition gives a zero-operation assembly $Q_T$ of the current
discarded family inside $Z_T$, including when $D_T$ is empty.
The local split construction identifies
\begin{equation}
  M_T\sqcup Z_T\xrightarrow{\cong} C_T.
  \label{eq:local-split}
\end{equation}
On a survivor region $U_{T,i}$, this map is $a_i$ composed with the inverse
of the survivor identification; on $Z_T$ it is the inclusion into $C_T$.
The survivor regions are open and closed, since connected components of a
manifold are open and closed.  They partition $M_T$, and the discarded regions
partition $Z_T$.  The maps therefore paste smoothly, with smooth inverses
defined on the partition $V_{T,i}$ of $C_T$.  This proves the split
diffeomorphism.  The original local operation chain, from $C_T$ to
$M_{q_T}$, remains to be performed.

The substitution is a geometric operation.  If $R$ has pieces indexed by a
finite type $A$, append the pieces indexed by $D_T$ to form the disjoint union
indexed by $A\amalg D_T$.  Every connected-sum step in $R$ is lifted over the
extra carrier $Z_T$: its balls, punctured regions, collar,
and gluing map are included in the first summand, while the extra carrier is
left untouched in the second.  The lifted chain ends at
\(M_T\sqcup Z_T\).  Composing it with the local chain in
$S_T$ yields an assembly of $M_{q_T}$.

More explicitly, a step with sides $A,B$ and output $Y$ lifts to a step with
sides $A,B\sqcup Z_T$ and output $Y\sqcup Z_T$.  The second ball lies in
$B$, so its punctured complement is the old punctured $B$ together with all
of $Z_T$.  The two target regions are the old first region and the union of
the old second region with $Z_T$; the collar is the old collar included in
$Y\sqcup Z_T$.  Disjointness, coverage, smooth inverses and the two collar
gluing identities follow separately on these open summands.  Lifting each
step, and leaving a zero-length chain zero-length, gives the claimed chain.

To compose across the split diffeomorphism, distinguish the two possibilities
for the local chain $C_T\to M_{q_T}$.  If it has no step, transport the
assembled target directly across the split.  If it has a first step, pull
back that step's two-region decomposition by the split diffeomorphism and
keep its ball and gluing data; then append its remaining steps.  This changes
the source of an actual operation.  A diffeomorphism alone is not inserted as
a connected-sum step.

The index family after substitution is literally
\[
  A\amalg D_T,
\]
and its carrier at the left summand is the old tail carrier while its carrier
at the right summand is the current discarded piece.  The finite equivalence
\begin{equation}
  \operatorname{Tail}(T')\amalg D_T
       \;\cong\;\operatorname{Tail}(T)
  \label{eq:tail-equivalence}
\end{equation}
is bijective because every ledger event at least $T$ is either $T$ itself or is
at least $T'$ by successor minimality.  The two cases are disjoint, and every
index in the right-hand side belongs to one of them.  The equivalence preserves
the carrier pointwise, so reindexing does not alter any region or gluing map.

\section{Reverse induction}
\label{sec:finite-reverse-induction}

For $e\in\mathcal E$, let $\mathscr A_e$ denote an assembly of the family indexed by its tail
with target $M_{q_e}$.  We prove, by well-founded induction on the strict
reverse order of the finite ledger, that
\begin{equation}
  \forall e\in\mathcal E,\qquad
  \mathscr A_e\text{ exists}.
  \label{eq:reverse-invariant}
\end{equation}

At the terminal event $e=H$, the slice $M_H$ is empty.  Hence its local
conclusion has no survivor.  The tail is exactly the finite family $D_H$, and
the local reconstruction $S_H$ itself is an assembly of that family.  A
finite equivalence between the local indices and $\operatorname{Tail}(H)$
relabels this assembly without changing its carriers.

For an earlier event $e=T$, choose its immediate successor $T'$.  The
induction hypothesis supplies $\mathscr A_{T'}$.  Apply the target transport
along $d_{T,T'}^{-1}$, perform the substitution of
Section~\ref{sec:finite-substitution}, and then reindex by
\eqref{eq:tail-equivalence}.  The resulting object is $\mathscr A_T$.
The later whole-slice assembly is retained as one disjoint summand, which
can itself be disconnected, and all present discarded factors are added
exactly once.

No survivor is silently discarded.  Every survivor is represented by its
component transport to the next event, while every non-survivor appears in the
right summand of one substitution.  The tail equivalences prove both coverage
and non-duplication.

\begin{theorem}[Finite reverse-history reconstruction]
\label{prop:connected-sum-reconstruction}
\uses{lem:finite-event-references,lem:finite-assembly-transport,prop:v4-global-extinction}
\lean{PoincareMT.m72FiniteReconstruction}
\source{MT2007}\dcref{Corollary 15.4, pp. 358--359}
For every fixed flow and local reconstruction data as above, there are a ledger $\mathcal E$, a
finite type $K$, a bijection
\[
  \iota:K\xrightarrow{\cong}\mathcal D,
\]
and a family $P_k=P(\iota(k))$ with a smooth finite connected-sum assembly
\[
  \mathscr A:\{P_k\}_{k\in K}\longrightarrow M_0.
\]
The ledger is exactly \eqref{eq:finite-ledger}; every $P_k$ is a connected
non-survivor piece of the actual local conclusion at its event; and every
survivor piece carries the successor record
\eqref{eq:successor-slab}--\eqref{eq:successor-component}.  The target carrier
is nonempty and connected.
\end{theorem}

\begin{proof}
The ledger is finite by local finiteness and nonempty because it contains the
extinction event.  Choose its first event $e_0$.  Equation~\eqref{eq:reverse-invariant}
gives an assembly of $\operatorname{Tail}(e_0)$ with target $M_{q_{e_0}}$.
The first-event anchor is zero: if it were a surgery time, it would be
nonnegative and strictly less than $e_0\leq H$, hence would belong to
$\mathcal E$, contrary to minimality.
The inverse predecessor transport is therefore a diffeomorphism
$M_{q_{e_0}}\to M_0$.  Transport the assembly along it and use the equivalence
$\operatorname{Tail}(e_0)\cong\mathcal D$.  The resulting family is indexed by
the exact global non-survivor type.  Enumerating this finite type by a finite
ordinal supplies $K$ and $\iota$.

The initial identification in the global certificate is a diffeomorphism from
$M$ onto $M_0$.  Its image transports the assumed connectedness and the
nonemptiness of $M$ to $M_0$.  The successor choices constructed above are
stored separately in the conclusion and their target times are shown to lie
in the same ledger by the interval bounds in \eqref{eq:finite-ledger}.
\end{proof}

\section{The event-indexed sphere-factor ledger}
\label{sec:m82-factor-ledger}

The next chapter classifies the exact family just constructed.  To make the
interface explicit, suppose in addition that $M$ is simply connected and that
Theorem~\ref{thm:exact-sphere-factors} is applied to this same reconstruction conclusion.
For each enumerated index $k\in K$, it supplies
\[
  \operatorname{kind}(P_k)=\mathrm{positive\ spaceform},
  \qquad \sigma_k:P_k\xrightarrow{\cong}S^3.
\]
Thus the event-indexed ledger is the tuple
\begin{equation}
  (K,\iota,(P_k)_{k\in K},(\sigma_k)_{k\in K},\mathscr A),
  \label{eq:m82-ledger}
\end{equation}
where $\iota$ is the very bijection supplied by the reconstruction.  Its
assembly target is still the same $M_0$; no new target map or independently
chosen factor list is introduced.

The tuple~\eqref{eq:m82-ledger} records exactly what the
downstream reduction needs.  Neither this ledger nor the reconstruction claims
prime-decomposition uniqueness.

\section{The reconstruction in the proof}
\label{sec:finite-history-contracts}

The reconstruction uses one flow with local finiteness, its
actual surgery times and the initial identification, together with
an empty slice at an actual event time of that same flow.  Its local
topological input consists, at every ledger event, of the
connected local pieces, survivor regions, cap correspondence, the empty-event
no-survivor rule, and the collared assembly of the incoming slice.  The
successor proof additionally needs regular-slab diffeomorphisms on intervals
with no surgery.

The next chapter receives the exact family $P_k$, its piecewise
kind data and the assembly $\mathscr A$, and constructs the maps $\sigma_k$.
It must not
replace them by a separately classified list or a new assembly.  The finite-history proof
does not prove the sphere-bundle obstruction, positive-spaceform uniformization,
or the identity $S^3\mathbin{\#}X\cong X$; those are downstream obligations.

Simple connectedness was not used in the reverse induction.  It enters
the next step, where the actual discarded factors are proved to be
three-spheres and the finite assembly is reduced to a sphere.
